\chapter{Probleme und Herausforderungen}
\label{cha:probleme}
Im Projektverlauf traten technische und organisatorische Herausforderungen auf, etwa bei Energieoptimierung, Hardwareintegration und Testbedingungen. Diese Punkte werden hier strukturiert festgehalten.

Die Dokumentation der Probleme hilft, Entscheidungen im Nachhinein zu verstehen und ähnliche Risiken in zukünftigen Projekten frühzeitig zu erkennen.

\section*{PCB-Aufwand deutlich höher als geplant}

\textbf{Problem:} Im Projektplan waren 9 Stunden für das PCB eingeplant. Tatsächlich wurden ca.\ 56 Stunden aufgewendet.

\textbf{Grund:} Das PCB wurde von Grund auf in EasyEDA Pro neu entwickelt, da kein geeignetes Open-Source-Design für den STM32WLE5JC mit allen benötigten Schnittstellen existiert. EasyEDA ist ein browserbasiertes PCB-Design-Tool, das eng mit LCSC und JLCPCB verknüpft ist und dadurch Bauteilauswahl und Fertigung vereinfacht. Hinzu kamen mehrere Iterationen für RF-Design, Antennenmatching und Komponentenplatzierung. Nach Erhalt der gefertigten Platine wurden außerdem Bugs gefunden und behoben, unter anderem ein fehlerhafter TPS63900-Spannungsregler. Das fertige Design wird auf OSHWLab als Open-Source-Projekt veröffentlicht und nimmt am OSHWLab Open-Source-Contest teil.

\textbf{Konsequenz:} Der Mehraufwand beim PCB wurde durch Einsparungen in anderen Bereichen teilweise ausgeglichen; insgesamt liegt das Projekt im Stundenrahmen.

\section*{Fehler im TPS63900-Spannungsregler}

\textbf{Problem:} Nach Erhalt der gefertigten Platinen von JLCPCB wurde festgestellt, dass der TPS63900-Buck-Boost-Konverter für die 3,3-V-Systemversorgung nicht wie vorgesehen funktionierte. Der Chip arbeitete nicht korrekt, was eine Inbetriebnahme des Custom PCBs zunächst verhinderte.

\textbf{Ursache:} In Absprache mit einem Experten wurde die Wurzel des Fehlers im PCB-Design identifiziert. Die Ursache lag in einem Designfehler im TPS63900-Abschnitt des Schaltplans.

\textbf{Lösung:} Als physischer Workaround wurde der fehlerhafte TPS63900 vom LiPo-Akku-Plus und dem 3,3-V-Netz getrennt. Derselbe Chip wurde auf einem Breakout-Board verbaut und an die PCB-Header angesteckt. Das System läuft damit identisch wie vorgesehen, mit minimalem Overhead und gleicher Effizienz. Das korrigierte PCB-Design ist erarbeitet und für zukünftige Revisionen einsatzbereit.

\section*{Sensor-Wechsel}

\textbf{Situation:} Im Projektplan wurden bestimmte Sensoren als Kandidaten genannt. Im Verlauf wurden leistungsfähigere Alternativen gewählt:

\begin{itemize}
  \item BME280 wurde durch SHT45 (Temperatur, Feuchte) und BMP390 (Luftdruck) ersetzt, da diese Kombination genauer und energieeffizienter ist.
  \item VEML6075 wurde durch LTR390 ersetzt, da der LTR390 einen breiteren UV-Messbereich abdeckt und einen I2C-Ausgang statt eines analogen Ausgangs bietet.
\end{itemize}

\textbf{Bewertung:} Alle Sensoren laufen stabil. Die neuen Sensoren sind für den Außenbetrieb besser geeignet als die ursprünglich geplanten.

\section*{Lautstärkesensor entfallen}

\textbf{Problem:} Feature 7 (Sensor Umgebungslautstärke, MAX4466, 3 Stunden geplant) wurde aus dem Scope genommen.

\textbf{Grund:} Der Fokus lag auf den Kernfunktionen. Ein analoges MEMS-Mikrofon erfordert zusätzliche Signalverarbeitung, die im verfügbaren Zeitrahmen nicht sauber umgesetzt werden konnte.

\textbf{Lösung:} Das Feature wird als \glqq out of scope\grqq{} markiert. Die wichtigsten Umweltdaten (Temperatur, Luftfeuchtigkeit, Luftdruck, UV, Feinstaub, CO${}_2$) werden vollständig erfasst.

\section*{KI-Evaluation noch nicht abgeschlossen}

\textbf{Problem:} Für die KI-Evaluation waren 16 Stunden geplant. Bis zum Projektabschluss wurden erst ca.\ 2 Stunden investiert, da der PCB-Aufwand Priorität hatte.

\textbf{Lösung:} Die KI-Analyse ist bereits im HydroNode-Backend aktiv: Sensordaten werden automatisch auf Anomalien analysiert, auch mit den neuen Outdoor-Daten. Eine formale Evaluation mit Modellauswahl und Dokumentation ist als nächster Schritt geplant, sobald die Hardware vollständig stabil läuft.

\section*{Nicht geplanter Mehraufwand: Gehäuse nach Stevenson-Screen-Prinzip}

\textbf{Situation:} Das Gehäuse-Design und der 3D-Druck wurden im Projektplan nicht als eigenes Feature aufgeführt. Tatsächlich wurden ca.\ 12 Stunden für Design in Fusion 360 und mehrere Druckdurchläufe aufgewendet.

\textbf{Warum dieser Aufwand nötig war:} Für eine korrekte Messung von Aussentemperatur und Luftfeuchtigkeit darf sich das Gehäuse nicht durch Sonneneinstrahlung aufheizen. Deshalb wurde das Design an einen Stevenson-Screen angelehnt: übereinanderliegende Lamellen blockieren direkte Sonneneinstrahlung und Regen, lassen aber gleichzeitig Luft frei zirkulieren. Als Material wurde ASA gewählt, da ASA UV-beständig, witterungsfest und für den dauerhaften Außeneinsatz geeignet ist. Der Druck erfolgte auf einem eigens für das Projekt angeschafften 3D-Drucker.

\section*{UV-Sensorschutz: Acrylglas mit Korrekturfaktor}

\textbf{Problem:} Der UV-Sensor LTR390 muss geschützt verbaut werden und trotzdem UV-Licht empfangen können. Normales Glas oder transparentes PETG blockiert einen großen Teil der UV-Strahlung und verfälscht so die Messwerte. UV-transparentes Spezialglas (Quarzglas, Borosilikatglas) würde ca.\ 50 Euro kosten und war daher keine wirtschaftliche Option.

\textbf{Lösung:} Aus Kostengründen wurde auf Standard-Acrylglas zurückgegriffen. Eine ca.\ 2~mm starke Acrylglasplatte wurde mithilfe einer CNC-Maschine rund zugeschnitten, sodass sie passgenau in die Öffnung des Stevenson-Screen-Gehäuses eingesetzt werden kann. Da Acrylglas einen Teil der UV-Strahlung absorbiert, wurde ein empirischer Korrekturfaktor ermittelt: Messwerte mit und ohne Glasabdeckung unter direkter Sonneneinstrahlung ergaben ein Verhältnis von $6{,}8 / 5{,}2 = 17/13 \approx 1{,}31$, das als ganzzahlige Bruchrechnung direkt in der Firmware angewandt wird. Die korrigierten UVI-Werte werden im Payload übertragen.

\section*{Versionsverwaltung und Arbeitsablauf}

\textbf{Problem:} Git und GitHub sind seit mehreren Jahren Teil meiner täglichen Arbeit. Ich nutze Pull Requests und GitHub Actions routiniert, was die Zusammenarbeit und die Integrationsprozesse wesentlich erleichtert. Im Verlauf dieses Projekts ist mir jedoch ein wiederkehrendes organisatorisches Problem aufgefallen. Üblicherweise beginne ich mit einem Issue und erstelle davon einen Branch aus \texttt{main}, in dem ich die Umsetzung vorantreibe. Dabei kommt es häufig vor, dass zusätzliche Aufgaben, kleinere Bugfixes oder Ergänzungen der Dokumentation ebenfalls in denselben Branch einfließen. Dies führt zu Branches, die mehrere inhaltlich unterschiedliche Änderungen enthalten und dadurch die Nachvollziehbarkeit von Änderungen sowie das gezielte Review erschweren.

\textbf{Lösung:} Um die Codequalität und die Review-Prozesse zu verbessern, ist es ratsam, thematisch getrennte Arbeiten konsequent in eigenen Branches zu bearbeiten und für jede Änderung ein separates Issue und einen eigenen Pull Request anzulegen. Auf diese Weise bleiben Änderungen atomar, Revisionshistorie und Verantwortlichkeiten werden klarer und Merge-Konflikte lassen sich leichter adressieren.
